
\subsubsection{4.1 Training Dataset}

The APPS dataset (\texttt{codeparrot/apps} on HuggingFace) is chosen because it spans the full algorithmic difficulty range needed for RQ1 (introductory through competition) and provides executable test cases required for RLEF reward computation. Following test-case filtering, 4,449 problems are retained for full-scale training. Smoke-scale experiments use 500 problems.

\textbf{Coverage measurement (h-m1, Task C):} Reference solution coverage is computed by running each APPS reference solution against all included test cases using subprocess execution (5.0-second timeout, temporary directory isolation). A problem is counted as covered if at least one reference solution passes all test cases. Out of 361 competition-level problems, 308 are covered, yielding 85.32\% coverage.

\subsubsection{4.2 Baselines}

\textbf{SFT} is the primary baseline. It uses the same base model (DeepSeek-Coder-7B-base) and training data (APPS) as RLEF, differing only in the training objective. No other model architectures or pre-trained checkpoints are included, to avoid confounds from model variation.

\textbf{RLEF-Binary} is included for the reward formulation comparison (RQ3) but uses proxy estimates due to the infrastructure constraint described in Section 3.3.

\subsubsection{4.3 Implementation Notes}

\textbf{max\_new\_tokens=512:} This value is required to avoid truncation of Python solution outputs. The reward monitoring experiment (h-m2) used max\_new\_tokens=128, producing 0.0 non-zero reward across all difficulty buckets. Inspection indicates that 128 tokens is insufficient for complete Python function bodies on competition-level problems, causing mechanical test-case failure from truncated code. This parameter is therefore set to 512 for all training and evaluation in h-e1, h-m3, and h-m4.

\textbf{Tokenizer source:} The SFT checkpoint tokenizer has a class incompatibility in the experimental environment; \texttt{deepseek-ai/deepseek-coder-7b-base} is used as the tokenizer source for all downstream evaluation.

\textbf{Compute environment:} $5\times$ NVIDIA H100 NVL GPUs (95,830 MiB each). Smoke-scale SFT training: 63 steps, final loss 0.486. RLEF training: 62 GRPO steps completed; checkpoint save failed due to process termination by session timeout after the training loop exited. RLEF evaluation therefore uses proxy estimates.

\subsubsection{4.4 Evaluation Metrics}

\textbf{pass@1:} Fraction of benchmark problems solved on first generation (greedy decoding), evaluated via bigcode-evaluation-harness.

\textbf{$\Delta$ = pass@1(RLEF) - pass@1(SFT):} Performance gap at each difficulty level. A positive, increasing trend in $\Delta$ across difficulty levels is the primary observable prediction.

\textbf{JT z-score:} One-tailed z-statistic from the Jonckheere-Terpstra test, with p < 0.05 as the significance threshold for confirming an ordered trend.

